
\lesson{10/11}
[qualcosina\\
ending of the ideal classical gas]
$$pV=NkT$$
This solution is simple because the $Z$ factorizes (no interactions) apart from a
$N!$ factor.

For exercise, work out the entropy $S(T,V,N)\to S(E,V,N)$ and check it agrees with
the Sakur-Tatrode microcanonical entropy.
$$Z(\beta) = \frac{1}{N!}\left(\frac{V}{V_T}\right)^N=\frac{1}{N!}\left(
\frac{V(2\pi m)^{\flatfrac{3}{2}}}{h^3}\right)^N
\beta^{-\frac{3N}{2}}$$
so from this, the integral is
$$g(E) = \frac{1}{2\pi i}\int_{\gamma}^{} \dd\beta e^{\beta E}Z(\beta)=
\mathcal{N}\frac{1}{2\pi i}\int_{\gamma}^{} \dd \beta e^{\beta E}\beta^{
-\frac{3N}{2}} = \mathcal{N} I_N(E)
$$
To integrate this, Jordan's lemma must be used to close the contour $\gamma$:
for $E<0$ this is trivially zero, but for $E>0$ the residues
must be computed; for even $N$ the pole is in $\beta=0$, for odd $N$ the
result is non-trivial, since there is a branch cut in $Re(\beta)\in\left(-\infty,0
\right)$, due to the inverse square root.\\
The first case is
$$I_{N=2k}(E)=\frac{1}{2\pi i} = \text{Res}\eval{\left(e^\beta E\right)\beta^{-3k}
}_{\beta=0}$$
expanding in Laurent series yields the result
$$I_{N}(E) = \frac{E^{\frac{3N}{2}-1}}{\Gamma\left(\frac{3N}{2}\right)}$$
The other integral is non-trivial and yields the same result, so
$$I_N(E) = \vartheta(E)\frac{E^{\frac{3N}{2}-1}}{\Gamma\left(\frac{3N}{2}
\right)}$$
and the density turns out to be
$$g(E)=\frac{1}{N!}\left[\frac{V(2\pi mE)^{\flatfrac{3}{2}}}{h^3}\right]^N
\frac{1}{\Gamma\left(\frac{3N}{2}\right)}\frac{1}{E}$$
which is the same result found from the microcanonical ensemble!

The probability of finding a microstate at energy $E$ is
$$P(E) = g(E)e^{-\beta E}\propto e^{-\beta E +\frac{3N}{2}\ln E}$$
which is maximal when the exponent is maximum:
$$0=\pdv[]{}{E}\left(-\beta E + \frac{3N}{2}\ln E\right)=-\beta+\frac{3N}{2E}=0
\iff E = \frac{3N}{2}\frac{1}{\beta} = \frac{3}{2}NkT
$$
which is also the mean value!!

From the from of $p(E)$ it is clear that the entropic and thermodynamic terms
``compete" and make the function peak at $\langle E\rangle$.
\subsubsection{Relation to Maxwell-Boltzmann}
[boh integrale in dd mu di rho etc, 5 minuti]
\section{Equipartition of energy}
Consider a system of $N$ \textit{distinguishable} particles with $6N$ canonical
coordinates, indexed as
$$\left(q^I,p_I\right)\text{, with }I=1,2...N\cdot d.$$
In general,
$$Z(\beta,V,N) = \int_{}^{} \underbrace{\frac{\dd^{dN}q\dd^{dN}p}{h^{dn}}}_{
:=\dd\mu}e^{-\beta H(q,p)}=\int_{}^{} \dd \mu e^{-\beta H}$$
the ``useful'' Hamiltonians are those which are quadratic in the $Nd$ momenta and
in the $Nf_q$ positions, where $f_q\leq d$ is the number of \textit{non-cyclic}
coordinates.

\textbf{Example} For the ideal gas Hamiltonian, $d=3$ and $f_q=0$, since there
is no position dependency in $H$.\\
\textbf{Example} For a collection of $N$ $1-$dimensional harmonic oscillators,
$d=1$ and $f_q = 1$, since
$$H = \sum_{I}^{}\frac{p_i^2}{2m}+\frac{m}{2}\omega_i^2q_i^2$$

This class of Hamiltonians is useful \textit{because it contains these two cases}:
when the system is close to mechanical equilibrium, it resembles the harmonic
Hamiltonian, through the theory of small oscillations (cfr MMMC bachelor's
course).

For these Hamiltonians, the Euler relation holds, since $H$ is an eigenfunction
of the Euler operator:
$$\underbrace{\left(\sum_{I}^{}q^I\pdv{}{q^I} + \sum_{I}^{}p_I\pdv{}{p_I}\right)
}_\text{Euler operator}H=2H$$
Taking the statistical average of this relation yields
$$2\langle H\rangle =\left(\langle\sum_{I}q^I\pdv{H}{q^I}\rangle+
\langle\sum_{I}^{}p_I\pdv{H}{p_I}\rangle\right)$$
this average is taken with the M-B distribution: the momentum term is
$$\langle p_I\pdv{H}{p_I}\rangle=
\int_{}^{} \dd \mu \frac{1}{Z}\underbrace{\left(p_I\pdv{H}{p_I}\right)e^{-\beta H}
}_{:=-\frac{1}{\beta}\pdv{}{p_I}e^{-\beta H}} = 
\frac{1}{Z}\int_{}^{} \dd \mu -\frac{1}{\beta}p_I\pdv{}{p_I}e^{-\beta H}
$$
this can be integrated by parts and the boundary term is suppressed by a Gaussian
factor:
$$\langle p_I\pdv{H}{p_I}\rangle=
-\frac{1}{\beta}\frac{1}{Z}(-)\int_{}^{} \dd\mu e^{-\beta H}=
\frac{1}{Z}\frac{1}{\beta}Z = kT,\,\forall p_I$$
The position term is similar, but 
$$\pdv{H}{q^I}\neq0\text{ only for non-cyclical }q^I$$
and the computation is the exact same, since the boundary condition for a confined
system goes to zero $e^{-\beta H}\underset{q^I\to\infty}\longrightarrow 0$:
$$\langle \pdv{H}{q^I}q^I\rangle = kT$$
So in total
$$2\langle H\rangle = \left[\sum_{I=1}^{Nf_q}kT + \sum_{I=1}^{Nd}kt
\right] = kTN(d+f_q)$$
and the final result, the \textbf{equipartition of energy}, is
$$\boxed{\langle H\rangle = \frac{kT}{2}N(d+f_q)}$$

\subsection{Virial theorem}
A related result is the so called \textit{statistical Virial theorem}: consider
$f_q=d$ and assume that only the kinetic term depends on $p_I$:
$$\langle K\rangle = 2\langle \sum_{I}^{}\frac{p_I^2}{2m}\rangle =
2\frac{kT}{2}Nd.\footnote{did not understand the 2.}$$
The Hamilton equations tell that
$$\pdv{H}{q^I}=-\dot p_I =-F_I$$
The Virial of the system is defined as
$$\mathcal{V} = \sum_{i=1}^{N}\mathbf{q_i\cdot F_i} = \sum_{I}^{}q_IF_I$$
